\section{Problem 2}
\textbf{(a)}: 
\begin{problem}
    State the hypothesis of $J(q, r)$
\end{problem}
\begin{solution}
    $q + 2 \neq 0$ or $r^2 - 1 \neq 0$. 
\end{solution}

\newpage
\textbf{(b)}
\begin{problem}
    State the conclusion of $J(q, r)$
\end{problem}
\begin{solution}
    \begin{equation*}
        (q + 2)(r^2 - 1) \neq 0
    \end{equation*}
\end{solution}

\newpage
\textbf{(c)}
\begin{solution}
    \begin{equation*}
        L(q, r): \text{"If $\paren{q + 2}\paren{r^2 - 1} = 0$, then $q + 2 = 0$ and $r^2 - 1 = 0$."}
    \end{equation*}

    That is false. 
\end{solution}

\newpage
\textbf{(d)}
\begin{solution}
    \begin{equation*}
        K(q, r): \text{"If $\paren{q + 2}\paren{r^2 - 1} \neq 0$, then $q + 2 \neq 0$ or $r^2 - 1 \neq 0$."}
    \end{equation*}

    That is true.
\end{solution}
  
\newpage
\textbf{(e)}
\begin{solution}
    \begin{equation*}
        \forall r \in \R, \exists q \in \R, \bracks{(q + 2 \neq 0 ) \vee \paren{r^2 - 1 \neq 0} \implies \paren{(q + 2)(r^2 - 1) \neq 0}}
    \end{equation*}

    This statement is true. 
\end{solution}

(f): This statement is true. 